% Gerado de docs/monografia/00-pre-textuais.md por scripts/monografia_latex.py.
% Não edite aqui: edite o .md e gere de novo.

Ponteiros estão entre os pontos mais difíceis do ensino de programação em C, e parte da dificuldade não é falta de conhecimento, mas concepções errôneas: ideias estáveis e coerentes para quem as tem, que resistem à explicação correta. Uma resposta errada, sozinha, não distingue quem não sabe de quem sabe errado, e os dois casos pedem intervenções diferentes. Este trabalho apresenta o Structura, um sistema de ensino de ponteiros em C que diagnostica cada premissa do conteúdo cruzando três sinais que o estudante fornece numa mesma tela (a resposta, a razão que a sustenta e a certeza declarada) e conduz cada caso por um percurso próprio: a quem não sabe, uma escada de passos curtos com retorno imediato; a quem sustenta uma concepção errônea, uma demonstração em que ele prevê o que um programa faz e vê o que a máquina fez; a quem erra depois da escada, uma revisão com o código animado linha a linha. O conteúdo cobre 13 tópicos, com 79 premissas, um catálogo de 61 concepções errôneas com procedência declarada, 402 passos, 80 demonstrações e 79 cenas animadas. A verificação foi feita sem depender da opinião de quem escreveu o material: um conferidor compila e executa cada trecho de código e confere o que o texto afirma, inclusive o estado que cada cena mostra, e dez perfis de estudante simulado atravessam todos os tópicos para conferir o comportamento descrito. O trabalho não avalia o sistema com estudantes; o estudo com participantes está desenhado e documentado como trabalho futuro. Conclui-se que é possível construir um sistema desse tipo com conteúdo e comportamento conferíveis, e que a verificação mecânica e a leitura humana encontram defeitos de natureza diferente e se completam.

\vspace{\onelineskip}
\noindent
\textbf{Palavras-chave:} ponteiros em C; concepções errôneas; diagnóstico; ensino de programação; verificação por compilador.
